\documentclass[11pt,a4paper]{article}

\usepackage[utf8]{inputenc}
\usepackage[T1]{fontenc}
\usepackage[ngerman]{babel}
\usepackage{amsmath,amssymb,amsthm}
\usepackage{geometry}
\usepackage{hyperref}
\usepackage{setspace}

\geometry{margin=2.5cm}
\onehalfspacing

\title{\textbf{Das Bamberg-Modell}\\
Eine arithmetische Grundlage für Teilchen, Felder und Konstanten}
\author{Thomas Hoffbauer}
\date{\today}

\begin{document}
\maketitle

\begin{abstract}
Wir stellen das Bamberg-Modell als eine vollständig interne,
endliche und arithmetische Fundamentstruktur vor,
aus der Teilchen, Wechselwirkungen,
Massenhierarchien und die Higgs-Struktur
ohne vorausgesetzte Quantenfeldtheorie emergieren.
Ausgehend von einer endlichen Zahlenkugel
mit modularem Familiensystem
werden Fermionen, Eichbosonen und das Higgs-Feld
als stabile Kohärenz- und Attraktorstrukturen identifiziert.
Physikalische Konstanten, insbesondere die Feinstrukturkonstante,
erscheinen als Zählverhältnisse und nicht als freie Parameter.
\end{abstract}

\section{Einleitung}

Das Standardmodell der Teilchenphysik beschreibt
die bekannten Elementarteilchen und ihre Wechselwirkungen
als lokale Quantenfeldtheorie
mit der Eichgruppe
$\mathrm{SU}(3)\times\mathrm{SU}(2)\times\mathrm{U}(1)$.
Trotz seines empirischen Erfolgs
bleibt die große Zahl freier Parameter,
die Herkunft der Massenhierarchien
sowie die Sonderrolle des Higgs-Feldes
konzeptionell unbefriedigend.

Das hier vorgestellte Bamberg-Modell
verfolgt einen alternativen Ansatz:
Nicht Felder oder Kontinua bilden den Ausgangspunkt,
sondern eine endliche arithmetische Struktur,
deren innere Kohärenz
die beobachteten physikalischen Strukturen
emergent hervorbringt.

\section{Grundstruktur der Bambergkugel}

\subsection{Endliche Zahlenbasis}

Die fundamentale Menge ist
\[
\mathcal{N}_{5000} = \{1,2,\dots,5000\}.
\]
Diese endliche Struktur ersetzt
kontinuierliche Raumzeitannahmen
durch eine vollständig diskrete Grundlage.

\subsection{Familiensystem}

Die Elemente werden nach ihren Restklassen modulo $12$
in vier Familien eingeteilt:
\[
\mathcal{F}=\{1,5,7,11\}\pmod{12}.
\]
Diese Wahl ist minimal, primzahlerhaltend
und symmetrisch.
Alle weiteren Strukturen des Modells
bauen auf dieser Einteilung auf.

\section{Kohärenz und Attraktoren}

\subsection{Morley--Walter-Kohärenz}

Zentrale Rolle spielt eine dreistellige
arithmetische Kohärenzrelation
zwischen Zahlentripeln.
Kohärente Tripel sind invariant
unter zyklischer Permutation
und bilden stabile lokale Orbits
in der Zahlenkugel.

\subsection{Attraktorbegriff}

Ein Attraktor ist eine Konfiguration,
die unter endlichen Permutationen
stabil bleibt.
Physikalisch beobachtbare Objekte
entsprechen genau solchen Attraktoren.

\section{Emergenz von Teilchen}

\subsection{Fermionen}

Fermionen sind stabile,
antisymmetrische Kohärenzorbits
mit fester Familienstruktur.
Das Pauli-Prinzip
ergibt sich als Unmöglichkeit
identischer Orbits.
Drei Generationen entstehen
als drei stabile Skalenfenster
innerhalb der endlichen Kugel.

\subsection{Eichbosonen}

Eichbosonen sind keine eigenständigen Objekte,
sondern Symmetriebewegungen
der Familienstruktur:
\begin{itemize}
\item Photon: globale Phasenrotation,
\item Gluonen: interne Permutationen der $A$-$B$-$C$-Familien,
\item $W^\pm$: Übergänge zwischen $E$ und $A,B,C$,
\item $Z$: normierter Mischmodus.
\end{itemize}

\section{Das Higgs-Feld}

\subsection{Definition}

Das Higgs-Feld ist keine zusätzliche Entität,
sondern die unvermeidliche
globale Familienimbalance:
\[
H =
(\rho_A,\rho_B,\rho_C,\rho_E)
- \frac14(1,1,1,1),
\]
wobei $\rho_X$
die relative Häufigkeit der Familie $X$ ist.

\subsection{Higgs-Mechanismus}

Die Massen der Fermionen
resultieren aus lokaler Verzerrung
dieser globalen Imbalance
in Verbindung mit der Kohärenzdichte.
Das Higgs-Boson erscheint
als niedrigster kollektiver Fluktuationsmodus
der gesamten Kugelstruktur.

\section{Physikalische Konstanten}

\subsection{Feinstrukturkonstante}

Die Feinstrukturkonstante wird definiert als
\[
\alpha
=
\frac{N_{\text{kohärent}}}{N_{\text{gesamt}}},
\]
wobei $N_{\text{kohärent}}$
die Zahl kohärenter elementarer Übergänge
und $N_{\text{gesamt}}$
die Zahl aller möglichen Übergänge ist.
Für $N=5000$ ergibt sich
\[
\alpha^{-1}\approx 137.
\]
Der Wert ist damit eine Zählgröße
und kein freier Parameter.

\section{Diskussion}

Das Bamberg-Modell
erklärt Teilchen, Felder und Konstanten
aus einer einheitlichen,
endlichen Struktur.
Es benötigt weder
eine fundamentale Raumzeit
noch zusätzliche Freiheitsgrade.
Das Standardmodell erscheint
als effektive Beschreibung
einer tieferliegenden arithmetischen Ordnung.

\section{Schlussfolgerung}

Physik wird im Bamberg-Modell
nicht auf Felder gegründet,
sondern auf stabile Zahlensymmetrien.
Teilchen sind Attraktoren,
Wechselwirkungen Bewegungsmodi,
und das Higgs-Feld
eine strukturelle Notwendigkeit
endlicher Kohärenz.

\begin{thebibliography}{9}

\bibitem{Gaillard}
M.~K.~Gaillard,
\emph{The Standard Model of Particle Physics},
arXiv:hep-ph/9812285.

\bibitem{Novaes}
S.~F.~Novaes,
\emph{Standard Model: An Introduction},
arXiv:hep-ph/0001283.

\bibitem{Dittmaier}
S.~Dittmaier, M.~Schumacher,
\emph{The Higgs Boson in the Standard Model},
arXiv:1211.4828.

\bibitem{Weinberg}
S.~Weinberg,
\emph{The Quantum Theory of Fields},
Cambridge University Press (1995).

\end{thebibliography}

\end{document}
